\documentclass[a4paper, serif, 10pt]{moderncv}

%% moderncv
% color and style
\moderncvcolor{black}
\moderncvstyle{banking}

%% page layout/margins
\usepackage[top=2.5cm, bottom=2.5cm, left=1.5cm, right=1.5cm]{geometry}

\nopagenumbers{}

%% Babel config for Norwegian language
% ref:
% - https://latex3.github.io/babel/guides/locale-norwegian.html
\usepackage[norsk]{babel}

%% fontspec
\usepackage{fontspec}

\IfFontExistsTF{Linux Libertine}{
  \setmainfont[Ligatures={TeX, Common}, Numbers=OldStyle]{Linux Libertine}
}{}
\IfFontExistsTF{Linux Libertine O}{
  \setmainfont[Ligatures={TeX, Common}, Numbers=OldStyle]{Linux Libertine O}
}{}

%% CTeX
% CJK support, used to show CJK characters in the resume
%
% - heading=false: disable ctex heading
% - punct=kaiming: use kaiming punctuations styles for CJK
% - scheme=plain: use plain scheme, do not override \normalsize font size
% - space=auto: space settings for CJK characters
%
% CJK fonts are resolved per locale language, e.g. Noto Serif/Sans CJK JP
% for Japanese resumes, with fallback to the next region if not installed
%
% ref:
% - http://ctan.mirrorcatalogs.com/language/chinese/ctex/ctex.pdf
\usepackage[UTF8, heading=false, punct=kaiming, scheme=plain, space=auto]{ctex}

\IfFontExistsTF{Noto Serif CJK SC}{
  \setCJKmainfont{Noto Serif CJK SC}
}{}
\IfFontExistsTF{Noto Sans CJK SC}{
  \setCJKsansfont{Noto Sans CJK SC}
}{}

%% line spacing
\usepackage{setspace}
\setstretch{1.125}

%% URL styling - use normal text instead of monospace
\urlstyle{same}

%% Patch \httplink and \httpslink to handle full URLs
% The original moderncv macros always prepend http:// or https:// to URLs.
% This causes issues when the URL already has a protocol (e.g., https://example.com)
% resulting in malformed URLs like https://https://example.com.
% This patch checks if the URL already contains "://" and uses it directly if so.
% Note: We use \str_set:Nx (x-expansion) to expand macros like \@homepage first.
\ExplSyntaxOn
\str_new:N \g_yamlresume_url_str

\RenewDocumentCommand{\httpslink}{O{}m}{
  \str_set:Nx \g_yamlresume_url_str {#2}
  \str_if_in:NnTF \g_yamlresume_url_str {://}
    { \url{#2} }
    { \url{https://#2} }
}

\RenewDocumentCommand{\httplink}{O{}m}{
  \str_set:Nx \g_yamlresume_url_str {#2}
  \str_if_in:NnTF \g_yamlresume_url_str {://}
    { \url{#2} }
    { \url{http://#2} }
}
\ExplSyntaxOff

\name{Marius Ødegård}{}
\title{Dataingeniør med erfaring innen skyplattformer og datapipelines}
\phone[mobile]{+47 912 34 567}
\email{marius.odegard@example.no}
\homepage{https://mariusodegard.no}

\address{Storgata 12, Oslo, Oslo, Norge, 0150}{}{}

\extrainfo{{\small \faLinkedin}\ \href{https://www.linkedin.com/in/marius-odegard}{@marius-odegard} {} {} {} • {} {} {} 
{\small \faGithub}\ \href{https://github.com/mariusodegard}{@mariusodegard}}

\begin{document}

\maketitle

\section{Grunnleggende informasjon}

\cvline{}{Dataingeniør med seks års erfaring fra å bygge skalerbare datapipelines, datavarehus og sanntidsløsninger. Sterk bakgrunn i Python, SQL, dbt, Airflow og skytjenester, med fokus på datakvalitet, automatisering og kostnadseffektivitet.
%
\begin{itemize}
\item Utvikler og drifter satsvise og strømmende datapipelines for analyse og maskinlæring
\item Erfaren med modellering av datavarehus og datakontrakter i tverrfaglige team
\item Opptatt av robust dataobserverbarhet, dokumentasjon og kontinuerlig forbedring
\end{itemize}}
\section{Utdanning}

\cventry{aug. 2015–juni 2017}
        {Master, Informatikk, Poeng: 3,8/4,0}
        {Universitetet i Oslo}
        {\url{https://www.uio.no}}
        {}
        {\begin{itemize}
\item Spesialisering i datahåndtering og distribuerte systemer
\item Masteroppgave om strømmebehandling og datakvalitet i sanntid
\end{itemize}
\textbf{Kurs}: Distribuerte systemer, Databaser og datamodellering, Maskinlæring, Statistisk analyse, Skyarkitektur}

\cventry{aug. 2012–juni 2015}
        {Bachelor, Informatikk, Poeng: 3,6/4,0}
        {Norges teknisk-naturvitenskapelige universitet}
        {\url{https://www.ntnu.no}}
        {}
        {\begin{itemize}
\item Grunnleggende og avansert programmering
\item Prosjektarbeid i tverrfaglige team med fokus på data og systemutvikling
\end{itemize}
\textbf{Kurs}: Programmering i Python, Algoritmer og datastrukturer, Operativsystemer, Nettverk, Databasesystemer}
\section{Arbeidserfaring}

\cventry{mars 2021–Nå}
        {Senior dataingeniør}
        {Nordlys Data AS}
        {\url{https://www.nordlysdata.no}}
        {}
        {\begin{itemize}
\item Ledet utviklingen av en skybasert dataplattform på Azure for 40+ interne brukere
\item Designet og implementert satsvise og strømmende datapipelines med Python, Databricks og Airflow
\item Etablerte datakontrakter, datakvalitetstester og overvåking som reduserte dataincidenter med 45 \%
\item Veiledet juniorutviklere og innførte kodestandarder og CI/CD for datajobber
\end{itemize}
\textbf{Nøkkelord}: Azure, Databricks, Airflow, Python, Datakvalitet}

\cventry{aug. 2018–feb. 2021}
        {Dataingeniør}
        {Fjord Analyse AS}
        {\url{https://www.fjordanalyse.no}}
        {}
        {\begin{itemize}
\item Bygget datavarehus i Snowflake med dbt-modeller for salg, kunde og produkt
\item Automatiserte daglige ETL-jobber med Airflow og Python, noe som sparte 15 timer per uke
\item Samarbeidet med analyse- og produktteam om datakrav og selvbetjeningsløsninger
\item Forbedret dokumentasjon og datakatalog for kritiske datakilder
\end{itemize}
\textbf{Nøkkelord}: Snowflake, dbt, ETL, SQL, Python}

\cventry{aug. 2017–juli 2018}
        {Dataanalytiker}
        {Havblikk Energi AS}
        {\url{https://www.havblikkenergi.no}}
        {}
        {\begin{itemize}
\item Utviklet rapporter og dashbord for produksjon og vedlikehold
\item Renset og modellerte sensordata fra vindturbiner
\item Støttet datadrevne beslutninger med SQL, Python og Power BI
\end{itemize}
\textbf{Nøkkelord}: Power BI, SQL, Sensordata, Rapportering}
\section{Språk}

\cvline{Norsk}{Morsmål eller tospråklig nivå \hfill \textbf{Nøkkelord}: Norsk, Bokmål}
\cvline{Engelsk}{Fullt profesjonelt nivå \hfill \textbf{Nøkkelord}: C1, Teknisk dokumentasjon}
\section{Ferdigheter}

\cvline{Datapipelines}{Ekspert \hfill \textbf{Nøkkelord}: Python, SQL, Airflow, dbt, Spark}
\cvline{Skyløsninger}{Avansert \hfill \textbf{Nøkkelord}: Azure, AWS, Snowflake, Databricks, Docker}
\cvline{Datamodellering}{Avansert \hfill \textbf{Nøkkelord}: Dimensjonsmodellering, Datavarehus, Datakontrakter, Datakatalog}
\cvline{Dataobserverbarhet}{Viderekommen \hfill \textbf{Nøkkelord}: Great Expectations, Overvåking, Datakvalitet, Logging}
\section{Utmerkelser}

\cventry{nov. 2023}
        {Nordlys Data AS}
        {Årets dataprosjekt}
        {}
        {}
        {Tildelt for dataplattformen som ga betydelig raskere innsikt og lavere driftskostnader.}

\cventry{juni 2017}
        {Universitetet i Oslo}
        {Beste masteroppgave i datahåndtering}
        {}
        {}
        {Anerkjennelse for masteroppgaven om strømmebehandling og datakvalitet i sanntid.}
\section{Sertifikater}

\cventry{sep. 2022}
        {Microsoft}
        {Microsoft Certified: Azure Data Engineer Associate}
        {\url{https://learn.microsoft.com/en-us/certifications/azure-data-engineer/}}
        {}
        {}

\cventry{mai 2023}
        {dbt Labs}
        {dbt Analytics Engineering Certification}
        {\url{https://www.getdbt.com/certifications/}}
        {}
        {}
\section{Publikasjoner}

\cventry{nov. 2017}
        {Strømmebehandling med fokus på datakvalitet i sanntid}
        {Norsk informatikkonferanse}
        {\url{https://www.nik.no}}
        {}
        {\begin{itemize}
\item Presenterer en arkitektur for sanntidsdatakvalitet i strømmepipelines
\item Diskuterer praktiske erfaringer fra energisektoren
\end{itemize}}
\section{Prosjekter}

\cventry{jan. 2022–des. 2023}
        {Åpen kildekode-prosjekt for datakvalitet i strømmepipelines}
        {Sanntids datakvalitetsplattform}
        {\url{https://github.com/mariusodegard/realtime-dq}}
        {}
        {\begin{itemize}
\item Bygget med Python, Kafka og Great Expectations
\item Automatiserte tester og varsling for kritiske datastrømmer
\item Dokumenterte mønstre for datakontrakter og overvåking
\end{itemize}
\textbf{Nøkkelord}: Kafka, Python, Datakvalitet, Observerbarhet}

\cventry{apr. 2021–okt. 2021}
        {Intern portal for kuraterte datasett og dashbord}
        {Selvbetjent analyseportal}
        {\url{https://www.nordlysdata.no/cases/analyseportal}}
        {}
        {\begin{itemize}
\item Utviklet metadata-API og søk i datakatalog
\item Integrerte Power BI og Snowflake for enklere selvbetjening
\end{itemize}
\textbf{Nøkkelord}: Datakatalog, Power BI, API}
\section{Interesser}

\cvline{Friluftsliv}{Fjellturer, Langrenn, Kajakk}
\cvline{Teknologi}{Åpen kildekode, Dataplattformer, Skyarkitektur}
\section{Frivillig arbeid}

\cventry{jan. 2022–Nå}
        {Fagansvarlig for dataplattform}
        {Dataforeningen Oslo}
        {\url{https://www.dataforeningen.no}}
        {}
        {\begin{itemize}
\item Arrangerer fagkvelder og workshops om datapipelines og skyarkitektur
\item Bidrar til nettverksbygging mellom datafaglige miljøer i Oslo
\end{itemize}}
    
\end{document}